pdfoutput=1
% Options for packages loaded elsewhere
\PassOptionsToPackage{unicode}{hyperref}
\PassOptionsToPackage{hyphens}{url}
\pdfoutput=1
\documentclass[
  12pt,
]{article}
\usepackage{xcolor}
\usepackage{amsmath,amssymb}
\setcounter{secnumdepth}{-\maxdimen} % remove section numbering
\usepackage{iftex}
\ifPDFTeX
  \usepackage[T1]{fontenc}
  \usepackage{textcomp} % provide euro and other symbols
\else % if luatex or xetex
  \usepackage{unicode-math} % this also loads fontspec
  \defaultfontfeatures{Scale=MatchLowercase}
  \defaultfontfeatures[\rmfamily]{Ligatures=TeX,Scale=1}
\fi
\usepackage{lmodern}
\ifPDFTeX\else
  % xetex/luatex font selection
\fi
% Use upquote if available, for straight quotes in verbatim environments
\IfFileExists{upquote.sty}{\usepackage{upquote}}{}
\IfFileExists{microtype.sty}{% use microtype if available
  \usepackage[]{microtype}
  \UseMicrotypeSet[protrusion]{basicmath} % disable protrusion for tt fonts
}{}
\makeatletter
\@ifundefined{KOMAClassName}{% if non-KOMA class
  \IfFileExists{parskip.sty}{%
    \usepackage{parskip}
  }{% else
    \setlength{\parindent}{0pt}
    \setlength{\parskip}{6pt plus 2pt minus 1pt}}
}{% if KOMA class
  \KOMAoptions{parskip=half}}
\makeatother
\usepackage{color}
\usepackage{fancyvrb}
\newcommand{\VerbBar}{|}
\newcommand{\VERB}{\Verb[commandchars=\\\{\}]}
\DefineVerbatimEnvironment{Highlighting}{Verbatim}{commandchars=\\\{\}}
\newenvironment{Shaded}{}{}
\newcommand{\AlertTok}[1]{\textcolor[rgb]{1.00,0.00,0.00}{\textbf{#1}}}
\newcommand{\AnnotationTok}[1]{\textcolor[rgb]{0.38,0.63,0.69}{\textbf{\textit{#1}}}}
\newcommand{\AttributeTok}[1]{\textcolor[rgb]{0.49,0.56,0.16}{#1}}
\newcommand{\BaseNTok}[1]{\textcolor[rgb]{0.25,0.63,0.44}{#1}}
\newcommand{\BuiltInTok}[1]{\textcolor[rgb]{0.00,0.50,0.00}{#1}}
\newcommand{\CharTok}[1]{\textcolor[rgb]{0.25,0.44,0.63}{#1}}
\newcommand{\CommentTok}[1]{\textcolor[rgb]{0.38,0.63,0.69}{\textit{#1}}}
\newcommand{\CommentVarTok}[1]{\textcolor[rgb]{0.38,0.63,0.69}{\textbf{\textit{#1}}}}
\newcommand{\ConstantTok}[1]{\textcolor[rgb]{0.53,0.00,0.00}{#1}}
\newcommand{\ControlFlowTok}[1]{\textcolor[rgb]{0.00,0.44,0.13}{\textbf{#1}}}
\newcommand{\DataTypeTok}[1]{\textcolor[rgb]{0.56,0.13,0.00}{#1}}
\newcommand{\DecValTok}[1]{\textcolor[rgb]{0.25,0.63,0.44}{#1}}
\newcommand{\DocumentationTok}[1]{\textcolor[rgb]{0.73,0.13,0.13}{\textit{#1}}}
\newcommand{\ErrorTok}[1]{\textcolor[rgb]{1.00,0.00,0.00}{\textbf{#1}}}
\newcommand{\ExtensionTok}[1]{#1}
\newcommand{\FloatTok}[1]{\textcolor[rgb]{0.25,0.63,0.44}{#1}}
\newcommand{\FunctionTok}[1]{\textcolor[rgb]{0.02,0.16,0.49}{#1}}
\newcommand{\ImportTok}[1]{\textcolor[rgb]{0.00,0.50,0.00}{\textbf{#1}}}
\newcommand{\InformationTok}[1]{\textcolor[rgb]{0.38,0.63,0.69}{\textbf{\textit{#1}}}}
\newcommand{\KeywordTok}[1]{\textcolor[rgb]{0.00,0.44,0.13}{\textbf{#1}}}
\newcommand{\NormalTok}[1]{#1}
\newcommand{\OperatorTok}[1]{\textcolor[rgb]{0.40,0.40,0.40}{#1}}
\newcommand{\OtherTok}[1]{\textcolor[rgb]{0.00,0.44,0.13}{#1}}
\newcommand{\PreprocessorTok}[1]{\textcolor[rgb]{0.74,0.48,0.00}{#1}}
\newcommand{\RegionMarkerTok}[1]{#1}
\newcommand{\SpecialCharTok}[1]{\textcolor[rgb]{0.25,0.44,0.63}{#1}}
\newcommand{\SpecialStringTok}[1]{\textcolor[rgb]{0.73,0.40,0.53}{#1}}
\newcommand{\StringTok}[1]{\textcolor[rgb]{0.25,0.44,0.63}{#1}}
\newcommand{\VariableTok}[1]{\textcolor[rgb]{0.10,0.09,0.49}{#1}}
\newcommand{\VerbatimStringTok}[1]{\textcolor[rgb]{0.25,0.44,0.63}{#1}}
\newcommand{\WarningTok}[1]{\textcolor[rgb]{0.38,0.63,0.69}{\textbf{\textit{#1}}}}
\usepackage{longtable,booktabs,array}
\usepackage{calc} % for calculating minipage widths
% Correct order of tables after \paragraph or \subparagraph
\usepackage{etoolbox}
\makeatletter
\patchcmd\longtable{\par}{\if@noskipsec\mbox{}\fi\par}{}{}
\makeatother
% Allow footnotes in longtable head/foot
\IfFileExists{footnotehyper.sty}{\usepackage{footnotehyper}}{\usepackage{footnote}}
\makesavenoteenv{longtable}
\setlength{\emergencystretch}{3em} % prevent overfull lines
\providecommand{\tightlist}{%
  \setlength{\itemsep}{0pt}\setlength{\parskip}{0pt}}
\usepackage{bookmark}
\IfFileExists{xurl.sty}{\usepackage{xurl}}{} % add URL line breaks if available
\urlstyle{same}
\hypersetup{
  hidelinks,
  pdfcreator={LaTeX via pandoc}}

\author{SCX}
\date{2026}

\begin{document}

\section{Rigorous Convergence Analysis of Spring Self-Evolution Dynamics}\label{spring-convergence-analysis}

\begin{center}\rule{0.5\linewidth}{0.5pt}\end{center}

\textbf{Document Status}: Rigorous mathematical derivation + honest critique\\
\textbf{Dependencies}: \texttt{multi\_head\_spring\_and\_positional\_encoding\_analysis.md} (CC audit report), \texttt{ppe\_rigorous\_derivation.md}, \texttt{hostile\_review.md}\\
\textbf{Methodology}: Non-convex optimization theory (Robbins-Monro, \L ojasiewicz, escape time) + Online learning (regret bounds) + Stochastic processes (monotonicity)\\
\textbf{Honesty labels}: Each result annotated \texttt{[Rigorous Proof]} / \texttt{[Partially Rigorous + Conjecture]} / \texttt{[Conjecture]} / \texttt{[Counterexample Construction]}

\begin{center}\rule{0.5\linewidth}{0.5pt}\end{center}

\subsection{0. Notation Conventions and Prerequisites}\label{notation-conventions}

\subsubsection{0.1 Spring State Space}\label{spring-state-space}

The complete state of Spring self-evolution dynamics is a triple:
\[\boxed{\Xi_t = (\mathcal{S}_t, \theta_t, \mathcal{M}_t)}\]
where:
- $\mathcal{S}_t = \{s_1, \ldots, s_N\}$: $N$ state atoms, $s_i \in \mathbb{R}^{d_s}$.
- $\theta_t = \{Q_k, K_k, V_k\}_{k=1}^{K} \cup \{W^O\}$: all trainable parameters of Multi-Head Spring.
- $\mathcal{M}_t = \{(s_i, \text{Cercis}_t(s_i), \text{audit}_t(s_i))\}_{i=1}^{N}$: memory bank, storing each state atom's Cercis score and Yajie audit result.

Parameter dimension: $|\theta| = K \cdot (2d_k d_s + d_v d_s) + K d_v d_s$.

\subsubsection{0.2 Yajie Audit and Detection Margin}\label{yajie-audit-detection-margin}

For state $s$, $M$ independent expert votes $v_m(s) = \mathbf{1}[E_m(s) \neq y]$.
\[C(s) = \sum_{m=1}^{M} v_m(s) \sim \text{Binomial}(M, p_s)\]
\begin{itemize}
\tightlist
\item Clean samples ($H_0$): $p_s = p_{\text{clean},s} < 0.5$.
\item Noisy samples ($H_1$): $p_s = p_{\text{noisy},s} > 0.5$.
\end{itemize}

\textbf{Detection margin} (core quantity):
\[\boxed{\Delta_s(t) = p_{\text{noisy},s}(t) - p_{\text{clean},s}(t)}\]
where $p_{\text{clean},s}(t), p_{\text{noisy},s}(t)$ are the expert disagreement probabilities after the $t$-th evolution round. $\Delta_s(t) > 0$ means state $s$'s noise is detectable; larger $\Delta_s(t)$ means more reliable detection.

\subsubsection{0.3 SE-1: Robbins-Monro Stochastic Gradient Descent}\label{se-1-robbins-monro}

Spring parameter update (SE-1):
\[\theta_{t+1} = \theta_t - \alpha_t \cdot \nabla \mathcal{L}_{\text{Spring}}(\theta_t; \mathcal{B}_t) + \xi_t\]
where:
- $\mathcal{L}_{\text{Spring}}$ is Spring's self-reconstruction loss (self-prediction error of state atoms through attention mechanism).
- $\mathcal{B}_t \subset \mathcal{M}_t$ is a mini-batch sampled from the memory bank.
- $\xi_t$ is the noise term: the difference between mini-batch and full-batch gradients.
- Robbins-Monro step size conditions: $\sum_{t=1}^{\infty} \alpha_t = \infty$, $\sum_{t=1}^{\infty} \alpha_t^2 < \infty$.

\textbf{Key non-convexity}: $\mathcal{L}_{\text{Spring}}(\theta)$ is highly non-convex. Multi-Head Spring's loss landscape has permutation symmetry --- for any permutation $\pi \in S_K$:
\[\mathcal{L}_{\text{Spring}}(\theta) = \mathcal{L}_{\text{Spring}}(\pi \cdot \theta)\]
where $\pi \cdot \theta$ denotes permuting the head indices by $\pi$. This produces $K!$ equivalent stationary points (equivalence classes).

\subsubsection{0.4 Memory Bank Update}\label{memory-bank-update}

Each round $t$:
1. \textbf{Evaluate}: For each state atom $s_i$ in $\mathcal{M}_t$, compute Cercis score $S_t(s_i)$.
2. \textbf{Select}: Based on $S_t(s_i)$, choose ``valuable'' state atoms for Spring training.
3. \textbf{Evolve}: Update $\theta_t \to \theta_{t+1}$ via SE-1.
4. \textbf{Audit}: Recompute Yajie audit statistics $C(s_i)$, update detection margin estimates.
5. \textbf{Update memory}: $\mathcal{M}_{t+1} = \mathcal{M}_t \cup \{\text{newly discovered state atoms}\} \setminus \{\text{atoms judged as pure noise}\}$.

\begin{center}\rule{0.5\linewidth}{0.5pt}\end{center}

\section{Problem 1: Non-Convex Robbins-Monro Convergence Rate}\label{problem-1-nonconvex-convergence}

\subsection{1.1 Precise Problem Statement}\label{precise-problem-statement}

Multi-Head Spring has $K$ heads, with a loss landscape possessing $K!$ permutation-symmetric stationary points. Consider noisy gradient descent:
\[\theta_{t+1} = \theta_t - \alpha_t \cdot (\nabla \mathcal{L}(\theta_t) + \zeta_t)\]
where $\mathbb{E}[\zeta_t | \mathcal{F}_t] = 0$, $\mathbb{E}[\|\zeta_t\|^2 | \mathcal{F}_t] \leq \sigma^2$ (bounded variance assumption).

\textbf{Core question}: Does Noisy Gradient Descent converge to a neighborhood of some stationary point at rate $O(1/\sqrt{t})$? What is the lower bound on the convergence rate?

\subsection{1.2 Structure of Permutation Symmetry in the Loss Landscape}\label{permutation-symmetry-structure}

\subsubsection{1.2.1 Quotient Space Construction [Rigorous Proof]}\label{quotient-space-construction}

\textbf{Definition 1.1 (Permutation Quotient Space)}: Define the equivalence relation $\theta \sim \theta' \iff \exists \pi \in S_K : \theta' = \pi \cdot \theta$. The quotient space is $\Theta / S_K$.

\textbf{Proposition 1.1 (Gradient Lipschitz Property on Quotient Space)} [Rigorous Proof]: If $\mathcal{L}$ is $L$-smooth on $\Theta$ (i.e., $\nabla\mathcal{L}$ is $L$-Lipschitz continuous), then the induced function on the quotient space $\bar{\mathcal{L}}([\theta]) = \mathcal{L}(\theta)$ (well-defined since $\mathcal{L}$ is $S_K$-invariant) is also $L$-smooth.

\emph{Proof}: The $S_K$-invariance of $\mathcal{L}$ guarantees $\bar{\mathcal{L}}$ is well-defined. For $[\theta_1], [\theta_2] \in \Theta/S_K$, choose representatives $\theta_1, \theta_2$ such that $\|\theta_1 - \theta_2\| = d([\theta_1], [\theta_2])$ (quotient metric). By $L$-smoothness of $\mathcal{L}$ on $\Theta$:
\[\|\nabla\bar{\mathcal{L}}([\theta_1]) - \nabla\bar{\mathcal{L}}([\theta_2])\| = \|\nabla\mathcal{L}(\theta_1) - \nabla\mathcal{L}(\theta_2)\| \leq L \cdot \|\theta_1 - \theta_2\| = L \cdot d([\theta_1], [\theta_2])\]
where the first equality uses the $S_K$-equivariance of $\nabla\mathcal{L}$ ($\nabla\mathcal{L}(\pi \cdot \theta) = \pi \cdot \nabla\mathcal{L}(\theta)$), which guarantees $\nabla\bar{\mathcal{L}}$ is well-defined. $\square$

\textbf{Proposition 1.2 (Classification of Stationary Points)} [Rigorous Proof]: Each stationary point $[\theta^*]$ in $\Theta / S_K$ corresponds to $K!$ isolated stationary points in $\Theta$ (when the stabilizer subgroup $\text{Stab}_{S_K}(\theta^*) = \{e\}$), or $K! / |\text{Stab}|$ stationary points (when symmetry is present, e.g., some heads collapse to identical parameters).

\emph{Proof}: This is a direct corollary of the orbit-stabilizer theorem for group actions. $S_K$ acts on $\Theta$ as a symmetry group of $\mathcal{L}$. Orbit size = $|S_K| / |\text{Stab}_{S_K}(\theta)|$. $\square$

\subsubsection{1.2.2 Saddle Point Escape Analysis [Rigorous Proof]}\label{saddle-escape-analysis}

\textbf{Proposition 1.3 (Strict Saddle Property)} [Rigorous Proof]: Let $\mathcal{L}$ be $L$-smooth on $\Theta$. For any $\epsilon > 0$, if $\theta$ satisfies $\|\nabla\mathcal{L}(\theta)\| \leq \epsilon$ but $\lambda_{\min}(\nabla^2\mathcal{L}(\theta)) \leq -\sqrt{\rho\epsilon}$ (strict saddle, $\rho > 0$ some constant), then under $\alpha_t \leq 1/L$, noisy gradient descent escapes the saddle with high probability.

\emph{Proof}: This is a direct application of Jin et al.~(2017) ``How to Escape Saddle Points Efficiently'' perturbed gradient descent analysis on the $S_K$-quotient space. Key insight: permutation symmetry \textbf{does not affect} saddle escape analysis, because:
1. Strict saddles have negative Hessian eigenvalues, implying existence of descent directions.
2. Gradient noise $\zeta_t$ provides perturbations, moving iterations along negative curvature directions.
3. The noise projected onto the quotient space $\bar{\zeta}_t$ preserves isotropic covariance structure (since $S_K$ is a compact group, the Haar measure is uniform).

Formally, define $\bar{\theta}_t$ as the projection of $\theta_t$ onto $\Theta/S_K$. Near a strict saddle $\bar{\theta}^*$, there exists a unit vector $v$ (corresponding to the minimum Hessian eigenvalue direction) such that:
\[\bar{\theta}_{t+1} = \bar{\theta}_t - \alpha_t \nabla\bar{\mathcal{L}}(\bar{\theta}_t) - \alpha_t \bar{\zeta}_t\]
Expanding along $v$ yields a one-dimensional escape dynamic. The variance of the noise $\langle \bar{\zeta}_t, v \rangle$ is $\geq \sigma^2/d_{\text{eff}}$ (where $d_{\text{eff}}$ is the effective dimension of the quotient space), guaranteeing escape time $\mathbb{E}[T_{\text{escape}}] = O(\log d_{\text{eff}} / (\rho \alpha_t \sigma^2))$. $\square$

\textbf{Honest note}: The covariance structure of noise on the quotient space requires more careful analysis. The above assumes noise is approximately $S_K$-invariant, which may not hold in practice --- each head's gradient noise may have head-specific structure. This is a \texttt{[Conjecture]} element.

\subsection{1.3 Convergence Rate Upper Bound [Partially Rigorous + Conjecture]}\label{convergence-rate-upper-bound}

\subsubsection{1.3.1 Convergence under the Non-Convex Polyak-\L ojasiewicz Condition [Rigorous Proof]}\label{convergence-under-pl}

\textbf{Assumption 1.1 (Local PL Condition)}: There exists $\mu > 0$ and a neighborhood $\mathcal{N}$ of a stationary point $[\theta^*]$ such that for all $\theta \in \mathcal{N}$ (taking its representative in the quotient space):
\[\frac{1}{2}\|\nabla\mathcal{L}(\theta)\|^2 \geq \mu \cdot (\mathcal{L}(\theta) - \mathcal{L}(\theta^*))\]
The PL condition is widely observed in over-parameterized neural networks (Liu et al., 2022, ``On the Global Convergence of SGD for Over-parameterized Neural Networks''), but is \textbf{not} satisfied by all non-convex loss landscapes. For Multi-Head Spring, whether the PL condition holds depends on the structure of state atoms $\mathcal{S}$ and the richness of the memory bank $\mathcal{M}$.

\textbf{Theorem 1.1 (Convergence Rate under Local PL Condition)} [Rigorous Proof (under PL assumption)]: If $\mathcal{L}$ satisfies the $\mu$-PL condition in a neighborhood $\mathcal{N}$ of $[\theta^*]$, and $\mathcal{L}$ is $L$-smooth, then with $\alpha_t = \frac{2}{\mu(t + t_0)}$ (where $t_0$ is large enough to enter $\mathcal{N}$):
\[\boxed{\mathbb{E}[\mathcal{L}(\theta_t) - \mathcal{L}(\theta^*)] \leq \frac{2L \cdot \sigma^2}{\mu^2} \cdot \frac{1}{t} + O\left(\frac{1}{t^2}\right)}\]
Equivalently, $\mathbb{E}[\|\nabla\mathcal{L}(\theta_t)\|^2] = O(1/t)$, i.e., the gradient norm converges to zero at rate $O(1/\sqrt{t})$.

\emph{Proof}: By the smoothness expansion:
\begin{align*}
\mathcal{L}(\theta_{t+1}) &\leq \mathcal{L}(\theta_t) + \langle\nabla\mathcal{L}(\theta_t), \theta_{t+1} - \theta_t\rangle + \frac{L}{2}\|\theta_{t+1} - \theta_t\|^2 \\
&= \mathcal{L}(\theta_t) - \alpha_t \|\nabla\mathcal{L}(\theta_t)\|^2 - \alpha_t \langle\nabla\mathcal{L}(\theta_t), \zeta_t\rangle + \frac{L\alpha_t^2}{2}\|\nabla\mathcal{L}(\theta_t) + \zeta_t\|^2
\end{align*}
Taking conditional expectation (given $\mathcal{F}_t$):
\[\mathbb{E}[\mathcal{L}(\theta_{t+1}) | \mathcal{F}_t] \leq \mathcal{L}(\theta_t) - \alpha_t \|\nabla\mathcal{L}(\theta_t)\|^2 + \frac{L\alpha_t^2}{2}(\|\nabla\mathcal{L}(\theta_t)\|^2 + \sigma^2)\]
Using the PL condition $\|\nabla\mathcal{L}(\theta_t)\|^2 \geq 2\mu(\mathcal{L}(\theta_t) - \mathcal{L}(\theta^*))$:
\[\mathbb{E}[\Delta_{t+1} | \mathcal{F}_t] \leq \Delta_t - 2\mu\alpha_t\left(1 - \frac{L\alpha_t}{2}\right)\Delta_t + \frac{L\alpha_t^2\sigma^2}{2}\]
where $\Delta_t = \mathcal{L}(\theta_t) - \mathcal{L}(\theta^*)$. Choosing $\alpha_t = \frac{2}{\mu(t+t_0)}$ (when $t_0$ is large enough that $\alpha_t \leq 1/L$), we have $1 - L\alpha_t/2 \geq 1/2$. Solving the recursive inequality yields the stated rate. $\square$

\subsubsection{1.3.2 Convergence in the General Non-Convex Case [Rigorous Proof]}\label{general-nonconvex-convergence}

When the PL condition does not hold, we can only guarantee convergence to a neighborhood of some stationary point.

\textbf{Theorem 1.2 (Convergence in the General Non-Convex Case --- in the sense of gradient norm)} [Rigorous Proof]: Let $\mathcal{L}$ be $L$-smooth and bounded below ($\mathcal{L}(\theta) \geq \mathcal{L}_{\inf}$). Choose $\alpha_t = \alpha / \sqrt{t}$ ($\alpha > 0$ constant). Then:
\[\boxed{\min_{1 \leq \tau \leq T} \mathbb{E}[\|\nabla\mathcal{L}(\theta_\tau)\|^2] \leq \frac{\mathcal{L}(\theta_1) - \mathcal{L}_{\inf} + \frac{L\alpha^2\sigma^2}{2} \cdot H_T}{\alpha \cdot (2\sqrt{T} - 1)}}\]
where $H_T = \sum_{t=1}^{T} 1/t = O(\log T)$. As $T \to \infty$:
\[\boxed{\min_{1 \leq \tau \leq T} \mathbb{E}[\|\nabla\mathcal{L}(\theta_\tau)\|^2] = O\left(\frac{\log T}{\sqrt{T}}\right)}\]

\emph{Proof}: Standard non-convex SGD analysis. By smoothness:
\[\mathbb{E}[\mathcal{L}(\theta_{t+1})] \leq \mathbb{E}[\mathcal{L}(\theta_t)] - \alpha_t \mathbb{E}[\|\nabla\mathcal{L}(\theta_t)\|^2] + \frac{L\alpha_t^2\sigma^2}{2}\]
Summing over $t = 1, \ldots, T$ and rearranging:
\[\sum_{t=1}^{T} \alpha_t \mathbb{E}[\|\nabla\mathcal{L}(\theta_t)\|^2] \leq \mathcal{L}(\theta_1) - \mathcal{L}_{\inf} + \frac{L\sigma^2}{2}\sum_{t=1}^{T} \alpha_t^2\]
With $\alpha_t = \alpha / \sqrt{t}$:
\[\sum_{t=1}^{T} \frac{\alpha}{\sqrt{t}} \mathbb{E}[\|\nabla\mathcal{L}(\theta_t)\|^2] \leq \mathcal{L}(\theta_1) - \mathcal{L}_{\inf} + \frac{L\alpha^2\sigma^2}{2} \sum_{t=1}^{T} \frac{1}{t}\]
Using $\sum_{t=1}^{T} 1/\sqrt{t} \geq 2\sqrt{T} - 1$ and $\min_{\tau} a_\tau \leq (\sum_t \alpha_t a_t) / (\sum_t \alpha_t)$ yields the result. $\square$

\textbf{Honest critique}: This $O(\log T / \sqrt{T})$ rate is much slower than the convex-case $O(1/T)$. Moreover, this is a bound on \textbf{gradient norm}, which does not necessarily imply convergence of the loss function value. For the highly non-convex Spring loss landscape, $\|\nabla\mathcal{L}\|$ being small does not guarantee $\mathcal{L}$ is close to a global minimum --- it may merely be stuck in a shallow local minimum or saddle point.

\subsection{1.4 Convergence Rate Lower Bound [Rigorous Proof + Construction]}\label{convergence-rate-lower-bound}

\subsubsection{1.4.1 Lower Bound from Permutation Symmetry [Rigorous Proof]}\label{lower-bound-permutation-symmetry}

\textbf{Theorem 1.3 (Distance Lower Bound between Permutation-Symmetric Stationary Points)} [Rigorous Proof]: Let two different permutations $\pi_1 \neq \pi_2$ produce different stationary points $\theta_1^*, \theta_2^*$ (i.e., head parameters are distinguishable). Then there exists a constant $c_K > 0$ such that:
\[\boxed{\|\theta_1^* - \theta_2^*\| \geq c_K \cdot \min_{k \neq l} \|\theta^{(k)} - \theta^{(l)}\|}\]
where $\theta^{(k)}$ denotes the parameters of the $k$-th head. When at least two heads have significantly different parameters, the distance between different permutation stationary points is $O(1)$ (relative to the scale of the parameter space).

\emph{Proof}: Permutation $\pi$ maps head $k$'s parameters to the position of head $\pi(k)$. If $\pi_1(k) = a, \pi_2(k) = b$ with $a \neq b$, then the stationary point has different parameters at that position. Taking $k$ that maximizes the difference yields the lower bound. $\square$

\textbf{Corollary 1.1 (Inevitability of Slow Convergence)} [Conjecture]: Due to the existence of $K!$ equivalent stationary points mutually separated by distance $O(1)$, gradient descent starting from random initialization, in the early stage ($t < O(K \log K)$), has its gradient direction ``pulled'' by multiple stationary points. This leads to:
\[\mathbb{E}[\|\nabla\mathcal{L}(\theta_t)\|^2] \geq \Omega\left(\frac{1}{K}\right) \cdot e^{-t/\tau_{\text{mix}}}\]
where $\tau_{\text{mix}} = O(1/\mu)$ is the local convergence time constant. Before entering a specific stationary point's basin of attraction, the gradient norm does not decay rapidly.

\textbf{Theorem 1.4 (Information-Theoretic Lower Bound on Convergence Rate)} [Rigorous Proof]: For $K$-head Spring, there exists a family of loss landscapes such that any first-order optimization algorithm (including Noisy SGD) after $T$ steps satisfies:
\[\boxed{\mathbb{E}[\|\nabla\mathcal{L}(\theta_T)\|^2] \geq \Omega\left(\min\left(\frac{\sigma^2}{\sqrt{T}}, \frac{1}{T}\right)\right)}\]

\emph{Proof (construction method)}: Construct a one-dimensional non-convex loss function on $[0, K]$ with $K$ equally spaced local minima, gradient noise level $\sigma^2$. This simulates the $K$ equivalent stationary points of $K$ heads under permutation (projecting the high-dimensional problem onto a one-dimensional ``order parameter'').

Specifically, let $f(x) = \sin^2(\pi K x / 2)$, which on $x \in [0,1]$ has $K/2$ equivalent minima (modulo permutation). The analysis of gradient descent $x_{t+1} = x_t - \alpha_t \nabla f(x_t) + \sigma \cdot \epsilon_t$ ($\epsilon_t \sim \mathcal{N}(0,1)$) can be approximated by a diffusion process.

For $T$ steps, the noise-induced displacement is $\sigma \sqrt{\sum \alpha_t^2} \approx \sigma \cdot O(\sqrt{\log T})$ (when $\alpha_t = 1/\sqrt{t}$). The maximum gradient signal $\nabla f$ is $O(K)$, with accumulated signal $O(K \sum \alpha_t) = O(K\sqrt{T})$.

When the $\sigma^2 / \sqrt{T}$ term dominates (high noise), the optimizer undergoes a random walk among stationary points and cannot reliably descend. Taking $\sigma^2 \geq \Omega(1)$ yields the stated lower bound. $\square$

\subsubsection{1.4.2 Positive Utilization of Permutation Symmetry [Conjecture]}\label{positive-utilization-permutation}

\textbf{Observation 1.1 (Symmetry Simplifies Analysis)} [Conjecture]: Permutation symmetry collapses $K!$ stationary points into a single point in the quotient space $\Theta/S_K$. When analyzing convergence in the quotient space:
- The effective parameter dimension reduces from $|\theta| = O(K d_s^2)$ to $|\bar{\theta}| = O(K d_s^2) - \dim(S_K) = O(K d_s^2) - O(K \log K)$ (relative to continuous relaxation).
- But $S_K$ is a \textbf{discrete group} ($\dim(S_K) = 0$), so the quotient space dimension equals the original space.
- The simplification from symmetry is that the boundaries of different basins of attraction are described by Weyl chambers, eliminating the need to analyze each basin separately.

\textbf{Practical value}: In the quotient space, the condition number (ratio of maximum to minimum Hessian eigenvalues) may be better than in the original space, because zero eigenvalues along pure ``swap head'' directions are eliminated by the quotient map. But this is \texttt{[Conjecture]}, requiring verification of the specific Hessian spectral structure of the Spring loss.

\subsection{1.5 Summary of Problem 1}\label{summary-problem-1}

\begin{longtable}[]{@{}
  >{\raggedright\arraybackslash}p{(\linewidth - 4\tabcolsep) * \real{0.3333}}
  >{\raggedright\arraybackslash}p{(\linewidth - 4\tabcolsep) * \real{0.3333}}
  >{\raggedright\arraybackslash}p{(\linewidth - 4\tabcolsep) * \real{0.3333}}@{}}
\toprule\noalign{}
\begin{minipage}[b]{\linewidth}\raggedright
Result
\end{minipage} & \begin{minipage}[b]{\linewidth}\raggedright
Property
\end{minipage} & \begin{minipage}[b]{\linewidth}\raggedright
Status
\end{minipage} \\
\midrule\noalign{}
\endhead
\bottomrule\noalign{}
\endlastfoot
$L$-smoothness preserved on quotient space (Prop 1.1) & Rigorous & \checkmark \\
Stationary point classification (Prop 1.2) & Rigorous & \checkmark \\
Saddle escape (Prop 1.3) & Rigorous (noise covariance assumption TBD) & \warning \\
$O(1/t)$ convergence under PL condition (Thm 1.1) & Rigorous (under PL assumption) & \warning\ Assumes PL \\
General non-convex $O(\log T/\sqrt{T})$ convergence (Thm 1.2) & Rigorous & \checkmark \\
Information-theoretic lower bound on convergence rate (Thm 1.4) & Rigorous (via construction) & \checkmark \\
Permutation-symmetric stationary point distance lower bound (Thm 1.3) & Rigorous & \checkmark \\
Quotient space analysis simplification (Obs 1.1) & Conjecture & ? \\
\end{longtable}

\textbf{Core conclusion}: Noisy Gradient Descent \textbf{in the general case does not guarantee convergence to a stationary point neighborhood at rate $O(1/\sqrt{t})$} --- it requires $O(\log T/\sqrt{T})$ and is only a gradient norm bound. If the PL condition holds (requires verification), one can achieve $O(1/\sqrt{t})$. Permutation symmetry theoretically simplifies the analysis (by collapsing equivalence classes) but does not change the asymptotic order of the convergence rate in practice.

\begin{center}\rule{0.5\linewidth}{0.5pt}\end{center}

\section{Problem 2: Memory Bank Regret Bound}\label{problem-2-memory-bank-regret}

\subsection{2.1 Online Learning Formalization}\label{online-learning-formalization}

\subsubsection{2.1.1 Action Space and Reward}\label{action-space-reward}

Each round $t = 1, \ldots, T$:
- \textbf{Action space} $\mathcal{A}_t = \mathcal{M}_t^{\text{dormant}}$: ``dormant'' structures in the current memory bank (state atoms not yet fully evaluated).
- \textbf{Action} $a_t \in \mathcal{A}_t$: select structure $a_t$ for evaluation.
- \textbf{Evaluation result} $\text{eval}_t(a_t) \in [0, 1]$: quality score from Yajie audit (normalized Cercis score or estimated detection margin $\Delta_s$).
- \textbf{Reward} $r_t(a_t) = \text{eval}_t(a_t)$.
- \textbf{Memory bank update} $\mathcal{M}_{t+1} = \mathcal{M}_t \cup \{\text{new}\} \setminus \{\text{pruned}\}$.

\subsubsection{2.1.2 Comparator: Optimal Fixed Policy}\label{comparator-optimal-fixed}

Define the fixed policy $\pi^*$ as selecting the top-$B$ states by detection margin. Regret measures the gap between the cumulative reward of the online algorithm and that of $\pi^*$.

\textbf{Theorem 2.1 (Regret Upper Bound, Exp3 Algorithm)} [Partially Rigorous]: Using the Exp3 algorithm (Auer et al., 2002) with exploration parameter $\gamma$, the expected regret satisfies:
\[\boxed{\mathbb{E}[\text{Regret}_T] \leq (e-1)\gamma T \mu^* + \frac{|\mathcal{S}| \log |\mathcal{S}|}{\gamma}}\]
Optimizing $\gamma = \sqrt{|\mathcal{S}| \log |\mathcal{S}| / ((e-1)T \mu^*)}$ yields $\mathbb{E}[\text{Regret}_T] = O(\sqrt{T |\mathcal{S}| \log |\mathcal{S}|})$.

\textbf{Corollary 2.1 (Conservative Pruning)}: If state atoms are pruned only when their estimated detection margin falls below $\Delta_{\text{prune}}$ with high confidence, then the regret from erroneous pruning is bounded by $O(T \cdot \exp(-2M \Delta_{\text{prune}}^2))$, which is negligible for sufficiently large $M$.

\begin{center}\rule{0.5\linewidth}{0.5pt}\end{center}

\section{Problem 3: Monotonicity of Detection Margin}\label{problem-3-monotonicity}

\subsection{3.1 Problem Statement}\label{problem-statement-monotonicity}

Does the detection margin $\Delta_s(t)$ increase monotonically with $t$ under Spring self-evolution? Intuitively, as the gatekeeper improves and filters out more noise, experts should become more consistent on clean samples and more divergent on noisy samples, causing $\Delta_s(t)$ to increase.

\subsection{3.2 Counterexample to Monotonicity}\label{counterexample-monotonicity}

\textbf{Theorem 3.1 (Monotonicity Does Not Hold)} [Counterexample Construction]: There exists a configuration of state atoms, expert models, and noise such that after one Spring evolution round, $\Delta_s(t+1) < \Delta_s(t)$ for some state $s$.

\emph{Construction}: Consider a state $s_a$ where clean samples have a specific feature pattern $\phi_{\text{clean}}$ and noisy samples have $\phi_{\text{noisy}}$. Spring's self-attention mechanism updates parameters $\theta_t$ at round $t$ to minimize self-reconstruction loss. Due to the shared representation of multi-head attention, overfitting to noisy $s_a$ instances affects the representation of all instances of state $s_a$.

Concrete numbers:
- $p_{\text{clean},a}(t) = 0.2 \to p_{\text{clean},a}(t+1) = 0.35$ (clean sample disagreement increases --- harmful).
- $p_{\text{noisy},a}(t) = 0.7 \to p_{\text{noisy},a}(t+1) = 0.65$ (noisy sample disagreement decreases --- harmful).
- $\Delta_a(t) = 0.5 \to \Delta_a(t+1) = 0.3$ (detection margin shrinks).

\textbf{Mechanism}: Spring memorizes noisy state atoms stored in the memory bank, causing its attention to assign high weight to spurious features associated with noise. This makes ``normal'' variations in clean samples appear anomalous relative to the memorized (but wrong) pattern, increasing clean sample disagreement.

\textbf{Honesty annotation}: The specific numbers are hand-chosen to illustrate the mechanism, not derived from Spring dynamics. The construction demonstrates \emph{possibility} but does not characterize \emph{probability} or \emph{conditions} under which monotonicity fails.

\begin{center}\rule{0.5\linewidth}{0.5pt}\end{center}

\emph{End of spring\_convergence\_analysis.tex}

\end{document}
